

\spis{Polując na neutrina, czyli z~pamiętnika fizyka. Dzień drugi}
{Joanna Zalipska}

\wtyt{\hspace*{-5.5cm}Polując na neutrina, czyli z~pamiętnika fizyka\\\hspace*{-5.5cm}Dzień drugi}
\waut{Joanna ZALIPSKA*}
\aafil{Narodowe Centrum Badań Jądrowych, Departament Badań Podstawowych, Zakład Fizyki Wielkich Energii}

Zaczynam kolejny dzień w~kopalni, w~której znajduje się detektor Kamiokande. Jak na razie obchód zrobiony i~nic się nie dzieje. Dane napływają, jest czas na pisanie. \marg{Artykuł ten jest kontynuacją serii opisującej obserwacje wykonywane za pomocą detektora Kamiokande w~Japonii. Pierwsza część ukazała się w~numerze~\href{https://www.deltami.edu.pl/2025/03/polujac-na-neutrina-czyli-z-pamietnika-fizyka-dzien-pierwszy/}{$\Delta\sb{25}^{3}$}.}


Zaczniemy od neutrin pochodzących ze Słońca. Oczywiście chcielibyśmy ich złapać jak najwięcej, żeby sprawdzić, czy
% strumień neutrin, czyli liczba neutrin, które do nas docierają, jest zgodna ze
liczba neutrin, które do nas docierają, jest zgodna z~liczbą neutrin przewidzianą w~standardowym modelu
%standardowym modelem
opisującym procesy zachodzące wewtąrz Słońca. 
Neutrina
pochodzące ze Słońca
mają różne energie. Nasz detektor niestety nie jest w~stanie zaobserwować ich wszystkich. Jeśli neutrino elektronowe ma zbyt niską energię, to wybity w~oddziaływaniu elektron ma za niski pęd, i~nie emituje promieniowania Czerenkowa albo emituje go za mało, żebyśmy byli w~stanie taki przypadek w~detektorze zrekonstruować. Próg na rejestrację neutrin słonecznych wynosi obecnie około 4~MeV, czyli badamy tylko neutrina elektronowe, które mają energię powyżej tej wartości. 
Problem w~tym, że aby mierzyć tak niskoenergetyczne neutrina, detektor musi być bardzo dokładnie dostrojony. Typowo tylko trzydzieści kilka fotopowielaczy spośród 11 tysięcy zarejestruje światło elektronu pochodzącego z~neutrina o~energii tak niskiej jak 4~MeV. Dlatego ważna jest precyzja pomiaru i~rekonstrukcji.
Tutaj z~pomocą przychodzą pomiary kalibracyjne wykonywane przy użyciu akceleratora elektronów, które pozwalają nam dostroić detektor.

\styt{Jak dostrajany jest detektor Super-Kamiokande?}

Rok temu brałam udział w~takiej kalibracji. Z~grupą około 10 osób, Japończyków, Amerykanów i~Europejczyków, montowaliśmy system pomiarowy i~zbieraliśmy dane. Jest to jedna z~trudniejszych kalibracji, które wykonuje się przy detektorze Super-Kamiokande.

\vspace*{\parskip}
\begin{dwieszpalty}
\vspace*{-6pt}
   \includegraphics[width=8.5cm]{Rysunek2-1.JPG}

    
    {\scriptsize Zdjęcie z~przygotowań do zbierania danych kalibracyjnych. Akcelerator znajduje się w~bunkrze po prawej stronie. Na zdjęciu widać szarą rurę, którą przemieszczają się elektrony. Zbiornik z~wodą jest poniżej -- naukowcy stoją na górnej pokrywie detektora Super-Kamiokande. Widoczna na zdjęciu zielona konstrukcja służy do przechowywania i~wprowadzania pionowych rur do zbiornika\vspace*{6pt}

    }\goodbreak
 
    Do kalibracji wykorzystujemy akcelerator elektronów o~energiach z~przedziału od 4~MeV do 18~MeV. Takie elektrony wpuszczamy do rury, w~której jest próżnia. Koniec rury umieszczony jest we wnętrzu zbiornika wodnego. W~ten sposób elektrony dosłownie wstrzeliwane są do wnętrza detektora, co imituje oddziaływanie neutrina. %Ponadto dokładnie znamy energię wstrzeliwanych elektronów (ponieważ sami je generujemy), dzięki czemu wiemy dokładnie, jakie elektrony powinniśmy zaobserwować, i~dzięki temu odpowiednio skalibrować detektor.
   Ponadto dokładnie znamy energię wstrzeliwanych elektronów (ponieważ sami je generujemy), co pozwala nam precyzyjnie przewidzieć, jakie elektrony powinny zostać zaobserwowane, a~tym samym odpowiednio skalibrować detektor.

\looseness-1{\spaceskip.33em minus.1emWygenerowane przez nas elektrony emitują promieniowanie Czerenkowa, które rejestrowane jest przez fotopowielacze na ściankach zbiornika wody. Następnie programy rekonstrukcyjne znajdują pierścienie promieniowania %(przykład na rys. 1)
pochodzące od elektronu oraz obliczają jego energię na podstawie ilości światła zarejestrowanego w~detektorze. Dzięki temu, że znamy początkową energię elektronów, możemy dostroić nasze programy rekonstrukcyjne tak, aby uzyskać poprawny wynik. }


\end{dwieszpalty}

    \styt{Zaginione neutrina słoneczne\dots?}
    
    W~wyniku wieloletnich badań wykazano, że neutrina mają pewną rzadką własność, a~mianowicie oscylują, czyli zmieniają w~czasie swoją tożsamość (na~przykład z~neutrin elektronowych na taonowe). \marg{Więcej o~oscylacjach neutrin można przeczytać w~\href{https://www.deltami.edu.pl/2024/10/oscylacje-neutrin/}{$\Delta\sb{24}^{10}$}.}
Odkrycia tego dokonano przy użyciu detektora Super-Kamiokande. W~czasie badań rejestrował on
\eject
\begin{dwieszpalty}
  tylko około 46$\%$ słonecznych neutrin elektronowych z~liczby przewidzianej przez standardowy model Słońca. Co oczywiście zaniepokoiło naukowców. W~poszukiwaniu odpowiedzi na pytanie, gdzie podziały się pozostałe neutrina elektronowe, przeprowadzono pomiary w~innym detektorze: Sudbury Neutrino Observatory (SNO) znajdującym się w~Kanadzie. Ten detektor jest w~stanie rejestrować oddziaływania neutrin mionowych i~taonowych ze Słońca. Okazało się, że brakujące neutrina elektronowe nie znikają, tylko zmieniają się w~neutrina taonowe albo neutrina mionowe, zanim jeszcze opuszczą Słońce. Za to odkrycie w~2015 roku część Nagrody Nobla z~fizyki przyznano Arthurowi McDonaldowi. 

\styt{Neutrina atmosferyczne i~ich przemiany}

Drugą część Nagrody Nobla w 2015 roku otrzymał Takaaki Kajita za odkrycie oscylacji neutrin atmosferycznych w~eksperymencie Super-Kamiokande. W~atmosferze ziemskiej w~wyniku oddziaływania promieniowania kosmicznego z~atomami tlenu i~azotu powstają neutrina mionowe i~elektronowe. Te neutrina są również rejestrowane przez eksperyment Super-Kamiokande. Łatwo jest je odróżnić od neutrin słonecznych, ponieważ mają znacznie wyższą energię – od kilkunastu GeV do TeV. W~detektorze rejestrujemy neutrina atmosferyczne przychodzące zarówno znad detektora -- po przejściu przez 1 kilometr góry, jak i~te, które powstały po drugiej stronie kuli ziemskiej i~przeniknęły przez całą średnicę Ziemi (oraz oczywiście neutrina pochodzące ze wszystkich innych stron wokół kuli ziemskiej). Zaobserwowano, że część neutrin mionowych, które przeszły przez kulę ziemską, znika po drodze.

Po tym odkryciu przeprowadzono pierwszy eksperyment wykorzystujący wiązkę neutrin mionowych (K2K) wyprodukowanych przy użyciu akceleratora (w~laboratorium KEK w~Japonii). Wiązkę tę wysyłano w~kierunku odległego o~250 km detektora Super-Kamiokande. Okazało się, że neutrina mionowe po drodze przemieniały się w~neutrina taonowe (oscylacje $\nu_{\mu} \rightarrow \nu_{\tau}$).

Do odkrycia pozostała jeszcze jedna transformacja -- neutrin mionowych w~neutrina elektronowe.
% Temu zadaniu sprostał obecnie działający eksperyment akceleratorowy T2K. Produkuje on intensywną wiązkę neutrin mionowych w~akceleratorze JPARC w~Japonii i~wysyła je do oddalonego o~295 km  detektora Super-Kamiokande. Eksperyment ten udowodnił pojawianie się neutrin
Zadanie to udało się zrealizować w~ramach eksperymentu akceleratorowego T2K, prowadzonego obecnie dzięki współpracy dwóch laboratoriów. Intensywna wiązka neutrin mionowych jest produkowana w~akceleratorze JPARC w~Japonii i~wysyłana do oddalonego o~295 km detektora Super-Kamiokande. W~eksperymencie T2K potwierdzono pojawianie się neutrin
elektronowych w~wiązce neutrin mionowych, czyli oscylacje $\nu_{\mu} \rightarrow \nu_e$. 

Na tym jeszcze nie koniec, ale\ldots\ zbliża się koniec dniówki w~kopalni. Dzisiaj nie było tak spokojnie, jak się można było spodziewać na początku szychty. W~pewnym momencie uruchomił się alarm wybuchu supernowej. Kiedy gwiazda umiera, zapada się grawitacyjnie i~wybucha. Wtedy emitowane jest zarówno światło, jak i~strumień niskoenergetycznych neutrin. Jak do tej pory, w~detektorze neutrin zaobserwowano takie zdarzenie tylko raz, w~1987 roku. Jak łatwo się domyślić, po uruchomieniu się alarmu nasza ekscytacja sięgnęła zenitu. Czyżbyśmy mieli szczęście i~w~trakcie dzisiejszej szychty wybuchła kolejna supernowa w~naszej Galaktyce? 

Niestety\ldots tym razem to tylko test. Poddano próbie naszą czujność, wywołując sztucznie alarm supernowej i~sprawdzając, czy zareagujemy zgodnie z~procedurami. Udało się, zaliczyliśmy test. I~pozostaje nam nadal oczekiwać -- na wybuch prawdziwej supernowej.
\end{dwieszpalty}

\input 09-tjz